% ======================================================================
% sec-intro.tex -- introduction. Lead instantiation chosen by \ifleaddatalog.
% ======================================================================
\section{Introduction}
\label{sec:intro}

Algebraic data provenance explains the outputs of a deterministic query---a
\emph{function} from an input database to an output
relation~\cite{green2007provenance,suciu2011probabilistic}. But many data
systems rest on a more ambiguous specification---a \emph{relation} that maps
an input scenario to several admissible outputs. A logic program with negation
may have several stable models; concurrent transactions may serialize in many
orders; distributed systems may reach different consistent decisions.
Classical data provenance cannot account for this ambiguity. It cannot even
pose its question---why, or why-not, a tuple appears---because the choice of
output is not yet made.

This paper develops an algebra of provenance over ambiguous specifications. We
introduce \emph{determination provenance}, which tracks the irrevocable
\emph{commitments} that resolve ambiguity. A \emph{determination} is a layered
sequence of commitments that narrows a relational specification until it
determines a single output---the semantic counterpart of what a correct
implementation does on one run. Classical provenance applies \emph{within}
each determination (recent work extends why/why-not provenance to a single
well-founded or stable model of a $\mbox{Datalog}^\neg$
program~\cite{bogaerts2026whynot}); determination provenance instead records
which commitments were needed and how query results depend on them \emph{across}
the admissible models.

\ifleaddatalog
\begin{example}[Why-not under ambiguity]
  \label{ex:intro}
  Take the program $r(c) \leftarrow \neg s(c);\ s(c) \leftarrow \neg r(c)$.
  It has two stable models: $\{r(c)\}$ and $\{s(c)\}$. Classical provenance
  cannot explain the absence of $s(c)$: in one model it is a why-not, in the
  other a why. The absence is not a blocked derivation but a \emph{semantic
  commitment}---the choice $\varphi_{r(c)=\mathbf{t}}$ that resolves the
  negative cycle. Determination provenance records exactly this: $r(c)$ holds
  under the determinations that commit it true, and that \emph{support} is the
  object provenance should annotate.
\end{example}
\else
\begin{example}[Why-not under ambiguity]
  \label{ex:intro}
  A relation $S(k,y)$ starts as $\{S(0,b)\}$; concurrent transactions
  $T_1\!:\textsf{ins }S(1,b)$, $T_2\!:\textsf{ins }S(2,d)$,
  $T_3\!:\textsf{del }S(2,d)$ run alongside a query
  $T_Q\!: Q(y)\drule S(k,y)$. Value $b$ appears regardless of ordering
  (\emph{robust}); value $d$ appears only when $T_Q$ observes $T_2$'s insert
  before an effective delete (\emph{contingent}). Before the ambiguity is
  resolved, the provenance question for $d$ is undetermined: under one
  resolution it is a why, under another a why-not.
\end{example}
\fi

\noindent
Each execution is one determination. But robustness and contingency are
properties of the \emph{space} of all determinations, not of any single one.
To reason about them algebraically we track each observable's \emph{support}:
the set of determinations under which it holds. Supports form the Boolean
shadow of a product provenance semiring, and the familiar power of algebraic
provenance transfers---robustness is full support; why/why-not questions are
read from the support; quantitative diagnosis from its measure.

Determination provenance carries more structure than this flat algebra.
Commitments are \emph{layered}: later ones may depend on earlier outcomes.
This layering induces a \emph{filtration}---a nested chain of sub-algebras
recording how much resolution each observable requires. A robust or impossible
observable has \emph{support depth}~$0$; a contingent one has positive depth,
naming the layer that first settles it. The key law is that provenance
composition cannot deepen the filtration: a positive relational-algebra query
produces no output deeper than its deepest input. This yields depth bounds,
compositional robustness, cross-specification comparison, and a connection to
probabilistic-database query complexity.

\paragraph{Contributions.}
\begin{enumerate}[label=(\roman*),nosep,leftmargin=*]
  \item \emph{A determination algebra} (Sections~\ref{sec:prelim}--\ref{sec:algebra}):
        a specification's resolutions form a complete determination space
        $\Dets$; provenance lifts to the product semiring $K^{\Dets}$ via an
        \emph{adequate consumer} whose sole law is that a nonzero annotation
        marks a true observable; supports form a Boolean semiring, and the
        canonical layer-prefix filtration measures support depth and is
        respected by provenance composition. Ordinary relational semiring
        provenance is recovered as the \emph{result-extensional} special case.
  \ifleaddatalog
  \item \emph{Lead instantiation, $\mbox{Datalog}^\neg$}
        (Section~\ref{sec:datalog}): sealing and choice commitments give a
        canonical basis of depth $k{+}d$; the classical negation semantics
        (stratified, well-founded, stable) are different reading depths of one
        shared filtration; sealing reproduces the support behavior of
        semirings with monus on stratified programs and extends to
        unstratified negation, where monus is undefined.
  \item \emph{Consequences} (Section~\ref{sec:discussion}): work regret vs.\
        semantic shift across specifications; a commitment-responsibility
        measure; and a bounded-treewidth algorithm for support and robustness.
        A second instantiation---transactional isolation---appears in
        Appendix~\ref{sec:transactions}.
  \else
  \item \emph{Lead instantiation, transactions} (Section~\ref{sec:transactions}):
        isolation levels (RC, SI, SER) are different resolving subsets of one
        shared filtration; determination depth is $\Theta(n)$ under scheduling
        discretion; SER and SI are incomparable in per-tuple support depth.
  \item \emph{Consequences} (Section~\ref{sec:discussion}): work regret vs.\
        semantic shift; a commitment-responsibility measure; and a
        bounded-treewidth algorithm for support and robustness. A second
        instantiation---$\mbox{Datalog}^\neg$---appears in
        Appendix~\ref{sec:datalog}.
  \fi
\end{enumerate}
